\chapter{Marco teórico y estado del arte}
\label{cap:marco}
Este capítulo reúne los conceptos necesarios para comprender las actividades de
Neuma: las habilidades musicales involucradas, la organización de la práctica,
el uso de juegos y la adaptación del desafío. Introduce también los fundamentos
del análisis vocal. Los procedimientos y decisiones concretas del prototipo se
desarrollan en los capítulos de diseño, implementación y pruebas. El estado del
arte complementa estos fundamentos con el análisis comparativo de antecedentes;
el procedimiento de la selección focalizada se describe en la sección~\ref{sec:metodo-revision}.

\section{Marco teórico}
\subsection{Conceptos musicales y habilidades de entrenamiento}
\label{sec:conceptos-habilidades}
\label{sec:altura-frecuencia}
\label{sec:intervalos-musicales}
\label{sec:afinacion-cents}
\label{sec:ritmo-pulso}
\label{sec:lectura-musical}
\label{sec:timbre-instrumentos}
\label{sec:escucha-conjuntos}
\label{sec:percepcion-produccion}
El entrenamiento auditivo exige distinguir las propiedades del sonido de las
acciones que se solicitan al usuario. La altura, o \textit{pitch}, es una
propiedad perceptiva; la frecuencia fundamental caracteriza la periodicidad
de la señal. Aunque suelen relacionarse, no son equivalentes, y los sonidos
vocales y musicales se apartan de la periodicidad ideal
\parencite[p.~1917]{decheveigne2002yin}. Esta diferencia importa al interpretar
una respuesta vocal a partir de una estimación computacional.

Un intervalo expresa la distancia entre dos alturas: es melódico si las notas
se suceden y armónico si suenan simultáneamente. Su denominación combina número
y cualidad, como en una tercera mayor
\parencite[sección~5.1]{hutchinsonIntervalos}. En el temperamento igual de doce
divisiones, la octava corresponde a una razón de frecuencias de dos a uno y se
divide en doce semitonos iguales. La frecuencia situada \(n\) semitonos respecto
de una referencia es
\begin{equation}
  f_n=f_{\mathrm{ref}}\,2^{n/12}.
  \label{eq:frecuencia-semitonos}
\end{equation}
La distancia entre dos frecuencias positivas puede expresarse en cents:
\begin{equation}
  c(f_2,f_1)=1200\log_2\left(\frac{f_2}{f_1}\right).
  \label{eq:distancia-cents}
\end{equation}
Una octava comprende 1200 cents y un semitono temperado, 100
\parencite{wolfeNotasFrecuencias}. Comparar una nota con una referencia externa
y comparar el intervalo entre dos notas producidas son tareas diferentes;
estas relaciones permiten expresar el objetivo y el desvío de cada respuesta.
En Neuma las utilizo para formular criterios de afinación y reproducción de
intervalos, cuyas tolerancias se definen en el diseño.

La dimensión temporal requiere distinguir pulso, tempo y ritmo. El pulso actúa
como referencia; el tempo indica su velocidad, expresable en pulsos por minuto
o BPM; el ritmo organiza sonidos y silencios respecto de esa referencia. El
compás agrupa los pulsos y las figuras representan duraciones relativas.
El puntillo añade la mitad del valor de una figura y la ligadura de prolongación
suma duraciones de notas de igual altura sin iniciar otra nota en la unión.
Las subdivisiones binarias, ternarias y los grupos de valoración especial
permiten describir patrones antes de fijar su velocidad
\parencite[secciones~4.2--4.5]{hutchinsonIntervalos}.
La sincronización sensoriomotora estudia la relación temporal entre estímulos
y acciones; la revisión de \textcite{repp2005synchronization} reúne investigaciones
con golpes de los dedos ante referencias auditivas. Este antecedente ofrece una
referencia para describir la ubicación temporal de las acciones en los ejercicios
rítmicos.

La lectura musical relaciona esa organización sonora con una representación
escrita. Las claves fijan referencias para interpretar las posiciones del
pentagrama, y la lectura de acordes exige reconocer relaciones entre notas.
En una inversión, la nota más grave puede diferir de la fundamental
\parencite[secciones~1.2, 6.1, 6.3 y 8.1]{hutchinsonIntervalos}. La primera vista
añade la ejecución inicial de material no practicado. El estudio de
\textcite{imai2021sightreading} sobre movimientos oculares en pianistas destaca
que los procedimientos difieren en la familiarización permitida antes de tocar.
Por ello, primera lectura, reconocimiento de notación y práctica repetida deben
describirse según la preparación y respuesta que admite cada tarea.

El timbre comprende atributos perceptivos no descritos únicamente por altura,
intensidad o duración. Su carácter multidimensional y su relación con la
identificación de fuentes incluyen aspectos como el ataque y la distribución
de energía entre frecuencias \parencite{mcadams2013timbre}. En conjuntos se añade
la organización perceptiva de fuentes simultáneas: \textcite{iverson1995streams}
encontró que atributos estáticos y dinámicos de timbres orquestales influían en
la segregación de secuencias y la identificación de melodías intercaladas.
La identidad instrumental debe distinguirse de funciones como melodía o
acompañamiento y de formas de realización como el arpegio; una misma armonía
admite distintas texturas \parencite[sección~14.3]{hutchinsonIntervalos}.

Finalmente, percibir una relación y producirla mediante la voz requieren
observaciones específicas. En participantes con formación musical escasa o nula
y sin formación vocal, \textcite{pfordresher2007singing} observaron que una
imitación vocal menos precisa no implicaba menor exactitud en discriminación
de alturas. Este resultado respalda distinguir percepción y producción dentro
de las tareas estudiadas. En Neuma, identificar un instrumento, leer una figura,
ejecutar un patrón temporal y cantar un intervalo constituyen respuestas
diferentes, aun cuando compartan contenido musical.

\subsection{Aprendizaje musical y organización de la práctica}
\label{sec:aprendizaje-practica}
\label{sec:aprendizaje-musical}
\label{sec:autorregulacion}
\label{sec:desempeno-aprendizaje}
\label{sec:practica-distribuida}
\label{sec:feedback-andamiaje}
\label{sec:referencias-melodicas}
La relación entre escuchar, leer y producir puede examinarse mediante el concepto
de audiación de Gordon: asimilación y comprensión mental de la música, diferenciada
de la percepción inmediata y de una imagen sonora sin comprensión musical
\parencite[pp.~3--4]{gordon2012sequences}. Tomo este concepto como referencia para
relacionar escucha, representación y producción en el diseño de las actividades:
qué relación musical presenta una consigna y qué necesita reconocer o producir
el usuario para resolverla.

La organización del estudio también forma parte de la práctica musical. El modelo
de \textcite[pp.~67--68]{zimmerman2002selfregulated} describe un ciclo de anticipación,
ejecución y autorreflexión: planificar objetivos, supervisar la actuación y revisar
los resultados para preparar el siguiente intento. En un estudio longitudinal
de práctica instrumental infantil, \textcite{mcpherson2001practice} observaron
diferencias individuales de autorregulación y un uso limitado de estrategias
de supervisión y corrección. Este antecedente sitúa la planificación y la revisión
del propio desempeño como aspectos relevantes de la práctica musical. Para Neuma,
retomo esa relación al organizar la elección de contenidos y la información
disponible después de un intento.

El resultado de una tarea constituye desempeño; el aprendizaje supone cambios
duraderos. La revisión de \textcite{soderstrom2015learning} distingue retención,
observada después de un intervalo temporal, de transferencia a condiciones o
tareas diferentes. La facilidad durante la práctica puede resultar insuficiente
para anticipar esos resultados. Repetir un ejercicio y resolver otro nuevo
permiten observar situaciones distintas. Por ello, vinculo el seguimiento del
desempeño con el contenido practicado y las condiciones del intento. Esta
distinción orienta tanto la variación de ejercicios como la interpretación
de los registros.

La distribución temporal ofrece otro criterio de organización. La práctica
distribuida separa encuentros con un contenido, mientras que la práctica masiva
concentra las repeticiones. El efecto de espaciamiento, o \textit{spacing effect},
es la ventaja de retención que puede producir esa distribución. Deben distinguirse
el intervalo entre episodios y el intervalo de retención hasta la prueba final;
\textcite{cepeda2008spacing}, en aprendizaje de hechos, mostraron que la separación
más favorable dependía del plazo de evaluación. La cantidad de práctica, las
separaciones y el momento de la prueba son, por tanto, variables necesarias
para estudiar ese efecto.

La evidencia musical introduce diferencias según el material. En aprendizaje de
canciones, \textcite{katz2021spacing} encontraron beneficios más claros para las
palabras que para las notas. En una melodía sin letra, \textcite{katz2025melodic}
observaron mejor recuerdo de notas e intervalos tras repasos espaciados, sin
ventaja global de una semana sobre dos días; el análisis se restringió a quienes
alcanzaron el criterio inicial. En identificación de seis intervalos,
\textcite{hua2015intervals} encontró un beneficio más pronunciado para intervalos
amplios y dependiente de las condiciones de presentación. Por su parte,
\textcite{molloy2012auditory} observaron mejoras entre sesiones y después del
entrenamiento en discriminación de frecuencias, y menor eficiencia de sesiones
extensas en fases posteriores; algunas comparaciones variaron también la cantidad
total de ensayos. Estos trabajos orientan mi propuesta de apoyar la continuidad
de la práctica y las oportunidades de repaso. La organización temporal debe
considerarse junto con el contenido y el volumen de trabajo, manteniendo abierta
su adecuación a cada actividad.

La retroalimentación permite revisar qué ocurrió y cómo continuar. El modelo
de \textcite{hattie2007feedback} diferencia información sobre tarea, proceso,
autorregulación y persona, relacionándola con las metas y los próximos pasos.
La noción de andamiaje de \textcite{wood1976tutoring} describe un apoyo que controla
componentes inicialmente fuera del alcance del aprendiz para que atienda a otros
que puede resolver. Aquel estudio se realizó con niños en una tarea de construcción.
Retomo el concepto para especificar qué parte del desafío resuelve una ayuda y
qué debe aportar el usuario, como ocurre al ofrecer dos notas de referencia,
solo una nota inicial o una consigna de producción libre.

En entrenamiento de intervalos, una melodía familiar puede actuar como referencia
asociativa. \textcite{smith1994familiar} observaron mejoras en identificación de
terceras y cuartas mediante etiquetas de melodías conocidas, con beneficio más
limitado en discriminación difícil. Su utilidad depende, por tanto, de la tarea.
Estos fundamentos permiten distinguir información sobre el error, una ayuda para
resolverlo y el resultado posterior. Las condiciones concretas de referencias,
metas y apoyos de Neuma se reúnen en el capítulo~\ref{cap:diseno}.

\subsection{Juegos serios, gamificación y motivación}
\label{sec:juegos-serios-gamificacion}
\label{sec:mecanicas-aprendizaje}
\label{sec:motivacion-autonomia}
\label{sec:metas-constancia}
\label{sec:evidencia-gamificacion}
Los juegos serios conservan la condición de juego y se orientan a propósitos que
exceden el entretenimiento. La gamificación utiliza elementos de diseño de juegos
en contextos que no constituyen por sí mismos un juego
\parencite{deterding2011gamification}. La distinción permite analizar tanto una
actividad completa como elementos que organizan la práctica. En Neuma distingo
las reglas y objetivos de cada juego de los apoyos transversales para organizar
el entrenamiento.

El aprendizaje basado en juegos vincula la actividad con resultados de aprendizaje
definidos. \textcite{plass2015foundations} plantean alinear las mecánicas con esos
objetivos y examinar los elementos que añaden procesamiento ajeno al contenido.
También describen la posibilidad de fallar con consecuencias reducidas y reintentar.
Utilizo estos criterios para revisar la correspondencia entre consigna, acción
y habilidad musical, y para organizar el tratamiento del error y el reintento.

La teoría de la autodeterminación aporta una perspectiva sobre la experiencia de
juego. \textcite{ryan2006motivation} consideran autonomía, competencia y vinculación
social. En sus estudios encontraron asociaciones entre autonomía y competencia
percibidas y disfrute o preferencia por continuar. Estas asociaciones ofrecen
criterios para analizar la motivación durante la práctica. De modo complementario,
\textcite{silverman2023streaks} estudiaron cómo representar una racha conservada o
interrumpida afecta la continuidad de conducta en contextos de consumo. En el
proyecto, considero las metas y rachas como apoyos a la participación, con registros
diferenciados del desempeño musical.

El metaanálisis de \textcite{sailer2020gamification} encontró efectos positivos
de gamificación educativa sobre resultados cognitivos, motivacionales y conductuales,
con heterogeneidad entre estudios. Los dos últimos fueron menos estables al
considerar trabajos con mayor rigor metodológico, y persistían límites para evaluar
su duración mediante pruebas demoradas. Esta evidencia fundamenta el interés
por integrar juego y práctica musical, atendiendo a las condiciones del diseño.
Para Neuma priorizo desafíos vinculados con la habilidad objetivo, superación
personal y autonomía con orientación; los criterios concretos se desarrollan
en el capítulo~\ref{cap:diseno}.

\subsection{Adaptación de dificultad y lógica difusa}
\label{sec:desafio-dda}
\label{sec:conjuntos-difusos}
El ajuste dinámico de dificultad, o DDA por \textit{Dynamic Difficulty Adjustment},
modifica condiciones del juego a partir de información de la interacción.
\textcite{hunicke2005dda} analiza requisitos de adaptación y presenta un sistema
que regula recursos en un juego de acción, destacando la coherencia de reglas y
devoluciones. Retomo ese criterio para relacionar los cambios de dificultad con
las reglas de cada actividad. Con Neuma busco sostener un desafío abordable a
partir del desempeño observado y conservar información que permita examinar
por qué se aplicó cada ajuste.

La lógica difusa permite expresar categorías con pertenencia gradual. Un conjunto
difuso \(A\), definido sobre un universo \(X\), asigna a cada elemento un grado
\(\mu_A(x)\) entre cero y uno; la pertenencia binaria es un caso particular
\parencite[pp.~338--340]{zadeh1965fuzzy}. Esta pertenencia describe compatibilidad
con una categoría y no debe interpretarse como probabilidad de dominio de una
habilidad. Las funciones que representan categorías como desempeño bajo o alto
son parte del modelo elegido.

Las reglas relacionan esas categorías con consecuencias. El experimento de
\textcite{mamdani1975linguistic} convierte reglas heurísticas de un operador en
control automático de una máquina de vapor. En la inferencia Mamdani, las
pertenencias activan reglas cuyas consecuencias difusas se combinan; la
desfusificación obtiene luego una salida numérica
\parencite{mathworksProcesoDifuso,mathworksDesfusificacion}. Este esquema permite
explicitar los criterios de ajuste y revisar la relación entre observaciones,
reglas y consecuencias. Esa posibilidad de inspección fundamenta su uso en el
proyecto. Los operadores, el centroide y las restricciones
de aplicación del proyecto se desarrollan en la sección~\ref{sec:diseno-dda}.

Adaptar el desafío y habilitar contenidos son decisiones diferentes: la primera
modifica condiciones de una actividad y la segunda regula acceso mediante reglas
de progresión. Su evaluación requiere identificar qué se observa, qué cambia y
cuándo se aplica. Separar ambas decisiones permite ajustar condiciones de práctica
conservando el progreso habilitado y comprobar cada mecanismo mediante sus
propios registros y casos de prueba.

\subsection{Audio digital y estimación de frecuencia fundamental}
\label{sec:audio-estimacion}
\label{sec:muestreo-audio}
El análisis vocal utiliza una representación digital de la señal. El muestreo
uniforme toma valores separados por \(1/f_s\) segundos, donde \(f_s\) es la frecuencia
de muestreo. Para señales limitadas en banda, la reconstrucción exige que el
contenido espectral quede por debajo de \(f_s/2\)
\parencite[apéndice~D, Sampling Theorem]{smithMuestreo}. Esa condición concierne
al espectro completo, no solamente a la frecuencia fundamental. El análisis de un
flujo utiliza segmentos finitos o ventanas, que pueden solaparse al incorporar
muestras nuevas.

YIN estima frecuencia fundamental mediante comparaciones temporales de la señal
consigo misma. Utiliza una función de diferencias, normalización por media
acumulada, búsqueda de un mínimo bajo umbral e interpolación del período
\parencite[pp.~1919--1921]{decheveigne2002yin}. Su fundamento es obtener un
desplazamiento compatible con la periodicidad. La implementación basada en YIN de
Neuma y sus diferencias frente al artículo se explican en la
sección~\ref{sec:estimacion-yin}. Esta base permite organizar la captura, el cálculo
y la comparación con el objetivo musical como etapas identificables del sistema.

Una estimación aceptada constituye un dato intermedio para evaluar la respuesta
musical. Deben distinguirse ausencia de señal suficiente, ausencia de periodicidad
admisible y una frecuencia válida que difiere del objetivo. Otras voces o una
referencia reproducida pueden entrar en el micrófono aunque se capture un único
canal. Por ello, delimito el análisis a una respuesta monofónica y defino
condiciones de captura coherentes con la tarea. La separación entre admisión
de señal y evaluación musical permite ofrecer una devolución acorde con el
estado observado.

La devolución también tiene una dimensión temporal. Android distingue latencia
de entrada, de salida y de ida y vuelta, que incorpora procesamiento
\parencite{androidLatenciaAudio}. El tiempo de cálculo de una ventana, el intervalo
entre estimaciones y la demora hasta una devolución visual describen magnitudes
diferentes. El diseño de las ventanas y el tratamiento de señal se presentan en
la sección~\ref{sec:implementacion-audio}; los criterios de evaluación técnica,
en la sección~\ref{sec:verificacion-audio}.

\subsection{Síntesis y relación con las actividades de Neuma}
\label{sec:sintesis-marco}
Los fundamentos anteriores permiten describir cada actividad por el contenido
musical que presenta, la respuesta que solicita y las condiciones en que la observa.
En el reconocimiento de instrumentos se identifica una fuente dentro de un banco
de estímulos y alternativas. La lectura de notas y acordes por nombre observa
reconocimiento de notación, mientras que la lectura rítmica con golpes en pantalla
registra inicios y relaciones temporales. Esta caracterización permite formular
un criterio de respuesta propio de cada actividad.

En los ejercicios vocales, reproducir una nota o un intervalo añade producción a
la escucha. La referencia ofrecida, las ayudas y el material condicionan la
interpretación del desempeño. Repetir el mismo ejercicio y resolver otro con
relaciones equivalentes permiten observaciones diferentes. La actividad de Simón
añade reproducción ordenada de una secuencia; su fundamentación específica sobre
memoria musical queda pendiente de completar.

La propuesta futura de escucha de conjuntos amplía el análisis hacia instrumentos
y funciones en grabaciones anotadas; su diseño se describe en el
capítulo~\ref{cap:diseno}, junto con las ayudas y los apoyos a la continuidad.
Distingo esas propuestas de las actividades implementadas para mantener la
trazabilidad entre fundamento, diseño y verificación.

En conjunto, adopto estos fundamentos para definir tareas musicales observables,
apoyos identificados, retroalimentación pertinente y reglas de adaptación
revisables. La evaluación técnica y las sesiones exploratorias previstas permiten
examinar el funcionamiento y el uso del sistema. La retención y la transferencia
duraderas requieren una evaluación pedagógica específica; las distingo de los
resultados de funcionamiento y participación.

\section{Estado del arte}
\label{sec:estado-arte}
Los antecedentes se comparan según la habilidad trabajada, la respuesta del
usuario, el mecanismo de adaptación y la evidencia disponible. Se reúnen
publicaciones sobre sistemas y métodos con documentación de aplicaciones de
entrenamiento. Esta última permite identificar prestaciones declaradas por sus
responsables, mientras que los estudios permiten examinar condiciones y resultados
de evaluación. La selección focalizada que sustenta esta versión se describe en
la sección~\ref{sec:metodo-revision}; no representa un inventario exhaustivo del área.

\subsection{Práctica musical digital y juegos educativos}
\label{sec:antecedentes-practica}
La combinación de escucha, lectura y producción ya aparece en herramientas de
entrenamiento musical. La documentación de Perfect Ear presenta ejercicios de
intervalos, escalas, acordes, canto, lectura y ritmo. EarMaster Cloud reúne
actividades auditivas y vocales con asignación de tareas y seguimiento de resultados
por docentes \parencite{perfectearFunciones,earmasterCloud}. Tomo estos repertorios
como referencia funcional para relacionar contenidos, modalidades de respuesta
e información de seguimiento. La documentación permite caracterizar la oferta
de práctica; los resultados educativos se examinan mediante los estudios
que presentan una evaluación.

Los juegos educativos permiten examinar cómo una tarea musical se incorpora a
reglas de interacción. iClef, de \textcite{barate2013iclef}, propone lectura de
partituras con cambios frecuentes de clave en dispositivos móviles. Los autores
informan uso en clases de estudiantes de música y valoraciones favorables de
diversión e implicación mediante cuestionarios. La evaluación de eficacia de
aprendizaje queda como una línea para ampliar. Para Neuma, el antecedente aporta
una relación explícita entre tarea de lectura y desafío lúdico: el contenido
musical forma parte de las acciones necesarias para avanzar.

En entrenamiento rítmico, \textcite{znidersic2025storytelling} compararon versiones
con y sin narrativa de \textit{Steady the Drums!}, un juego de realidad virtual
en el que patrones de percusión controlan unidades. El estudio de dos semanas
combinó cuestionarios, registros de juego y pruebas rítmicas iniciales y finales.
La narrativa recibió mejores valoraciones hedónicas, mientras que la versión sin
narrativa mostró mayor exactitud durante el juego. Hubo mejoras en medidas
rítmicas, pero no en todas las condiciones de tempo o dificultad. El diseño
compara dos versiones activas y no establece retención demorada ni eficacia frente
a ausencia de entrenamiento. Para Neuma, el antecedente muestra por qué conviene
separar disfrute, desempeño dentro del juego y resultado en una prueba musical.

En conjunto, estos trabajos ofrecen formas de vincular una habilidad con la
interacción. Comparo sus aportes atendiendo a la plataforma, la modalidad de
respuesta y la evaluación. Para Neuma, retomo como criterio que el desafío se
resuelva mediante la habilidad musical prevista y que los registros permitan
identificar la tarea efectivamente realizada.

\subsection{Análisis vocal y adaptación basada en desempeño}
\label{sec:antecedentes-audio-adaptacion}
El análisis vocal requiere distinguir la estimación acústica de la evaluación
musical. Sobre la base temporal de YIN, \textcite{mauch2014pyin} proponen pYIN:
generan candidatos con probabilidades a partir de una distribución de umbrales
y utilizan un modelo oculto de Markov para estimar una trayectoria de frecuencias.
\textcite{kim2018crepe} presentan CREPE, una red convolucional que opera sobre
la forma de onda y evalúan su precisión y robustez ante ruido. Ambos aportan
alternativas de procesamiento para examinar la estimación vocal. Su incorporación
al proyecto requiere considerar, además de precisión, demora y costo de ejecución
en el teléfono objetivo.

Organizo la selección técnica a partir de la implementación basada en YIN como
línea de base existente. Esta referencia permite plantear una comparación de
señales, errores, rechazos y tiempos bajo condiciones comunes antes de incorporar
otra alternativa. Así, la elección de DSP o aprendizaje automático se vincula
con las necesidades de la interacción vocal y las restricciones del dispositivo,
como se indica en el alcance.

La adaptación también cumple propósitos distintos. El CA-BAT de
\textcite{harrison2018cabat} utiliza teoría de respuesta al ítem, selección
adaptativa y generación automática para evaluar percepción del pulso. El
participante identifica cuál de dos fragmentos contiene una referencia alineada
con el pulso musical. Sus estudios apoyan el uso en laboratorio, con limitaciones
para evaluación en línea. El ajuste busca obtener información sobre la capacidad
del participante: es un antecedente de medición adaptativa, cuyo propósito difiere
de regular la experiencia de práctica mediante el DDA de Neuma.

SingPAD, presentado por \textcite{zhang2024singpad}, aporta registros de lectura
cantada obtenidos de SingMaster para estudiar seguimiento del conocimiento
o \textit{knowledge tracing}. Cada nota se vincula con conceptos musicales y
una etiqueta de respuesta derivada de evaluación automática. Los autores comparan
modelos para predecir respuestas a partir del historial. Su evaluación se concentra
en esa capacidad predictiva. Para Neuma, el aporte relevante es la vinculación
explícita entre contenido, producción vocal y registro, junto con la necesidad
de documentar cómo se obtienen las etiquetas de respuesta y qué datos se incluyen.

Estos antecedentes distinguen tres niveles: estimar la señal, interpretar una
respuesta y decidir la intervención siguiente. En Neuma relaciono esos niveles
mediante criterios musicales y reglas de adaptación explícitas. La lógica difusa,
presentada en la sección~\ref{sec:desafio-dda}, permite revisar cómo los indicadores
de desempeño activan una decisión. Ese vínculo orienta la estructura del motor
y de los registros con los que verifico su comportamiento.

\subsection{Comparación y aporte del proyecto}
\label{sec:comparacion-antecedentes}
La tabla~\ref{tab:antecedentes} resume el aporte de cada grupo de antecedentes
y el alcance de su comparación. Las diferencias entre objetos y preguntas de
evaluación permiten identificar qué retomo de cada línea de trabajo.

\begingroup
\small
\singlespacing
\setlength{\tabcolsep}{4pt}
\begin{longtable}{p{0.23\textwidth}p{0.31\textwidth}p{0.36\textwidth}}
\caption{Antecedentes y alcance de su relación con Neuma.}
\label{tab:antecedentes}\\
\toprule
Antecedente & Aporte relevante & Límite de comparación \\
\midrule
\endfirsthead
\toprule
Antecedente & Aporte relevante & Límite de comparación \\
\midrule
\endhead
\bottomrule
\endfoot
Perfect Ear y EarMaster Cloud & Repertorios de práctica; seguimiento docente en EarMaster Cloud. & Documentación funcional, sin evaluación independiente en esta comparación. \\
iClef & Juego móvil centrado en lectura de claves. & Aceptación en clases; eficacia de aprendizaje por ampliar. \\
Steady the Drums! & Comparación de narrativa, experiencia y ejecución rítmica. & Realidad virtual y dos versiones activas; sin prueba demorada de retención. \\
pYIN y CREPE & Alternativas de estimación de frecuencia fundamental. & Precisión de procesamiento; desempeño en el teléfono objetivo por medir. \\
CA-BAT & Selección adaptativa y evaluación de percepción del pulso. & Instrumento psicométrico; su adaptación no equivale a DDA de entrenamiento. \\
SingPAD & Datos de canto, conceptos musicales y predicción de respuestas. & Evaluación de modelos predictivos; no ensayo de eficacia de una intervención. \\
\end{longtable}
\endgroup

El aporte que propongo con Neuma es la integración y evaluación de un sistema
Android que articula tareas musicales, procesamiento local, retroalimentación
y adaptación inspeccionable. La relación entre estos componentes plantea
decisiones de ingeniería concretas: admitir una señal antes de evaluar su
afinación, conservar las condiciones musicales del intento y aplicar los ajustes
de dificultad en momentos compatibles con las reglas del juego.

Organizo esa integración de modo que las respuestas, las decisiones del motor
y el progreso puedan examinarse a partir de registros. Esto permite vincular
el criterio de cada actividad con su implementación y sus pruebas, y analizar
precisión del procesamiento, comportamiento del ajuste y experiencia de uso.
La contribución del trabajo se sostiene en esa construcción y en la evidencia
obtenida bajo condiciones explícitas. Los antecedentes proporcionan referencias
para justificarla y criterios para revisar sus resultados.
